\documentclass[12pt,a4paper]{article}
\usepackage[utf8]{inputenc}

\setcounter{page}{1}
\pagenumbering{arabic}
\usepackage{amsmath,amsfonts,amssymb,amsthm}
\usepackage{enumerate}
\usepackage[dvips]{graphicx,epsfig}
\usepackage{setspace}
\usepackage{geometry}
\usepackage{hyperref}
\usepackage{color}
\usepackage[british]{babel}
\usepackage[mmddyyyy]{datetime}
\renewcommand{\baselinestretch}{1.2}

\marginparwidth=0in \marginparsep=0in \oddsidemargin=0in \evensidemargin=0in \headheight=0in \headsep=0in \topmargin=0.5in \textheight=9in
\textwidth=6in
\setlength{\parindent}{0pt}

\begin{document}
\begin{tabular}{l r}
University of Essex  \hspace{35ex}& Session 2025-2026 \\
Department of Economics \hspace{35ex} & Autumn Term\\
Dr Ben Etheridge \hspace{35ex} &
\end{tabular}
\\
\begin{center}
\large{\textbf{EC466: Econometrics }\\
Problem Set \# 7: Staggered Difference-in-Differences}
\end{center}
\vspace*{0.5cm}

\textbf{Background:} Consider a policy intervention that was rolled out across different regions at different times. You have panel data from four regions (A, B, C, D) over six time periods (2018-2023), with the following average outcomes:

\vspace{1em}
\begin{center}
\begin{tabular}{lcccccc|c}
\hline
Region & 2018 & 2019 & 2020 & 2021 & 2022 & 2023 & Treatment Start \\
\hline
A & 45.2 & 46.1 & 52.8 & 55.3 & 57.9 & 60.2 & 2020 \\
B & 43.8 & 44.5 & 45.2 & 51.7 & 54.1 & 56.8 & 2021 \\
C & 44.6 & 45.3 & 46.1 & 46.8 & 53.4 & 55.9 & 2022 \\
D & 42.9 & 43.6 & 44.2 & 44.9 & 45.5 & 46.1 & Never \\
\hline
\end{tabular}
\end{center}
Each region has 100 observations per period, and treatment is irreversible once adopted.
\vspace{1em}

\begin{enumerate}
\item \textbf{Two-Way Fixed Effects Estimation}
\begin{enumerate}[(a)]
\item Write down the standard TWFE regression specification for this staggered adoption design. Define all variables clearly, including the treatment indicator $\mathcal{D}_{i,t}$.
\item Explain intuitively why a researcher might be tempted to use this TWFE specification. What causal parameter might they hope to estimate?
\item Calculate three $2\times2$ DiD estimates ``by hand'': (i) Region A vs.\ Region D comparing 2019 to 2020; (ii) Region B vs.\ Region D comparing 2020 to 2021; (iii) Region B vs.\ Region A comparing 2020 to 2021. Interpret each estimate.
\end{enumerate}

\item \textbf{The Goodman-Bacon Decomposition}
\begin{enumerate}[(a)]
\item According to the Goodman-Bacon (2021) decomposition, the TWFE estimator can be written as:
\begin{equation*}
\hat{\beta}^{DD} = \sum_{k \neq U} s_{k,U}\,\hat{\beta}^{DD}_{k,U} + \sum_{k \neq U}\sum_{\ell > k}\left[ s^{k}_{k\ell}\,\hat{\beta}^{DD,k}_{k\ell} + s^{\ell}_{k\ell}\,\hat{\beta}^{DD,\ell}_{k\ell} \right]
\end{equation*}
Explain what each component of this decomposition represents. What do the subscripts $k$, $U$, and $\ell$ denote?
\item Describe the three types of $2\times2$ comparisons that contribute to the TWFE estimator in a staggered design. Which of these comparisons is most likely to be problematic, and why?
\item The weights in the Goodman-Bacon decomposition depend on sample size and variance in the treatment variable. Explain why the weight $s_{k,U}$ for comparing early-treated group $k$ to the never-treated group $U$ includes the term $(n_k + n_U)^{2}$. Is this weighting scheme necessarily aligned with policy relevance?
\end{enumerate}

\item \textbf{Event Study Specifications}
\begin{enumerate}[(a)]
\item Write down a TWFE event study regression specification that includes relative time indicators (leads and lags). Explain why a researcher would use this specification instead of a simple TWFE model with a single treatment indicator.
\item Suppose you estimate a TWFE event study and find that the coefficients on the lead terms (pre-treatment periods) are statistically significant and economically large. Provide two distinct explanations for this finding: one that suggests a violation of identifying assumptions, and one that arises even when all assumptions hold.
\item Explain why ``binning'' of endpoints (e.g., using $\mathcal{D}^{<-K}_{i,t}$ for all periods more than $K$ periods before treatment) can contaminate the coefficients on other event-time indicators. What role does treatment effect heterogeneity play in this contamination?
\end{enumerate}

\end{enumerate}

\end{document}
